\begin{pastexam}[label={px:ncku-106-prog-02}]{成大資訊聯招 106 學年度 第 2 題}
（問答）

(10\%) Suppose there are 11 slots available from 250, please give the hashing function to insert
the keys, 42, 77, 85, 113, 315, 433, 474, 479, 574, 582, 698, to the 11 slots without collision by
using the division method, and show the hashing table.
\end{pastexam}

\begin{pastsoln}
「11 slots available from 250」沒有說清楚，這裡讀成：可用的位址是從 $250$ 開始的 $11$ 格，也就是
$250, 251, \ldots, 260$。 如果位址就是 $0..10$，把下面答案的 $250 +$ 拿掉即可，表的內容不變。

答案：$h(k) = 250 + (k \bmod 11)$。 雜湊表：
\begin{center}
\small
\begin{tabular}{c*{11}{c}}
\toprule
位址 & 250 & 251 & 252 & 253 & 254 & 255 & 256 & 257 & 258 & 259 & 260 \\
\midrule
鍵 & 77 & 474 & 574 & 113 & 433 & 698 & 479 & 315 & 85 & 42 & 582 \\
\bottomrule
\end{tabular}
\end{center}

\textbf{為什麼除數取 $11$。} 除法雜湊是 $k \bmod m$（資結書第~12.1~節），值落在 $0..m-1$。 要把 $11$ 個
鍵放進 $11$ 格又不碰撞，$11$ 個餘數必須兩兩不同：$m < 11$ 時餘數只有 $m$ 種，由鴿籠原理一定有兩個鍵相同；
$m > 11$ 時餘數可能超過 $10$，落到 $11$ 格以外。 所以取 $m = 11$，再加上起始位址 $250$。 $11$ 也是質數，
符合除法雜湊「$m$ 取質數」的建議。

\textbf{逐一計算 $k \bmod 11$}（寫成 $k = 11q + r$）：
\begin{center}
\small
\begin{tabular}{rlcc}
\toprule
$k$ & $11q + r$ & $r = k \bmod 11$ & $h(k)$ \\
\midrule
42 & $11 \cdot 3 + 9$ & 9 & 259 \\
77 & $11 \cdot 7 + 0$ & 0 & 250 \\
85 & $11 \cdot 7 + 8$ & 8 & 258 \\
113 & $11 \cdot 10 + 3$ & 3 & 253 \\
315 & $11 \cdot 28 + 7$ & 7 & 257 \\
433 & $11 \cdot 39 + 4$ & 4 & 254 \\
474 & $11 \cdot 43 + 1$ & 1 & 251 \\
479 & $11 \cdot 43 + 6$ & 6 & 256 \\
574 & $11 \cdot 52 + 2$ & 2 & 252 \\
582 & $11 \cdot 52 + 10$ & 10 & 260 \\
698 & $11 \cdot 63 + 5$ & 5 & 255 \\
\bottomrule
\end{tabular}
\end{center}
十一個餘數正好是 $0, 1, \ldots, 10$ 各一次，沒有碰撞，每一格都用到，載入因子是 $1$。
\end{pastsoln}
